
% =============================================================
\section{Branchpoint elimination}\label{app:branchpoint-elimination}
% =============================================================

Jaxley's internal representation introduces branchpoint pseudo-nodes
at cable junctions to enforce Kirchhoff's current law. Partitioning
$G$ into real ($R$) and branchpoint ($B$) blocks gives
\[
  \begin{pmatrix} G_{RR} & G_{RB} \\ G_{BR} & G_{BB} \end{pmatrix}
  \begin{pmatrix} \mathbf{v}_R \\ \mathbf{v}_B \end{pmatrix}
  =
  \begin{pmatrix} [G\mathbf{v}]_R \\ [G\mathbf{v}]_B \end{pmatrix}.
\]
Branchpoint nodes carry no membrane, so $c_b A_b = 0$ at $b$ and the
cable equation at $b$ reduces to the algebraic constraint
$[G\mathbf{v}]_B = 0$, giving $\mathbf{v}_B = -G_{BB}^{-1}G_{BR}\mathbf{v}_R$
and the Schur complement
\[
  \tilde{G} \;=\; G_{RR} - G_{RB}\,G_{BB}^{-1}\,G_{BR}.
\]
The same operator applies to the extracellular forcing:
$\tilde{G}\boldsymbol{\phi}_R$ gives the activating function at the
real compartments, where $\boldsymbol{\phi}_R$ is the field sampled at
real compartment centres only. The reduced operator acts only on real
compartments, so no $\boldsymbol{\phi}_B$ value enters the forcing. Our implementation delegates this assembly to a
private Jaxley helper rather than reconstructing it from scratch, so
the operator stays in sync with upstream solver changes.

